\documentclass[12pt]{article}
\usepackage[T1]{fontenc}
\usepackage[utf8]{inputenc}
\usepackage{lmodern}
\usepackage{textcomp}
\usepackage{enumitem}
\usepackage{geometry}
\usepackage{graphicx}
\usepackage{amsmath, amssymb}
\geometry{margin=1in}
\begin{document}
\begin{center}
Chemistry - Undergraduate Level Assessment 1 \
Total Marks: 40 \
Answer all questions.
\end{center}

\section*{Multiple Choice (10 marks)}
\begin{enumerate}[label=\arabic*.]
\item Which of the following is lower for argon than for neon?
\begin{enumerate}[label=(\alph*)]
    \item Melting point
    \item Boiling point
    \item Polarizability
    \item First ionization energy
\end{enumerate}
\item Calculate the equilibrium polarization of 13 C nuclei in a 20.0 T magnetic field at 300 K.
\begin{enumerate}[label=(\alph*)]
    \item 10.8 x 10^-5
    \item 4.11 x 10^-5
    \item 3.43 x 10^-5
    \item 1.71 x 10^-5
\end{enumerate}
\item Calculate the Q-factor for an X-band EPR cavity with a resonator bandwidth of 1.58 MHz.
\begin{enumerate}[label=(\alph*)]
    \item Q = 1012
    \item Q = 2012
    \item Q = 3012
    \item Q = 6012
\end{enumerate}
\item The number of allowed energy levels for the 55 Mn nuclide are:
\begin{enumerate}[label=(\alph*)]
    \item 3
    \item 5
    \item 8
    \item 4
\end{enumerate}
\item When the following equation is balanced, which of the following is true? \_\_ MnO 4- + \_\_ I- + \_\_ H+ \textless{}-\textgreater{} \_\_ Mn$_{2}$+ + \_\_ IO 3- + \_\_ H$_{2}$O
\begin{enumerate}[label=(\alph*)]
    \item The I- : IO 3- ratio is 3:1.
    \item The MnO 4- : I- ratio is 6:5.
    \item The MnO 4- : Mn$_{2}$+ ratio is 3:1.
    \item The H+ : I- ratio is 2:1.
\end{enumerate}
\end{enumerate}

\section*{True / False (10 marks)}
\textit{Answer True or False}
\begin{enumerate}[label=\arabic*.]
\setcounter{enumi}{5}
\item True or False: K+ has a smaller ionic radius than Sc$_{3}$+ because it is a smaller atom in the periodic table.
\item True or False: Infrared spectroscopy is a light-scattering technique that primarily analyzes molecular vibrations.
\item True or False: NH 3 is considered the strongest base in liquid ammonia because it has a high tendency to donate protons.
\item True or False: Silicon is most commonly found in nature as an oxide, primarily in the form of silica and silicate minerals.
\item True or False: Doubling the volume of a buffer solution by adding water significantly increases its pH.
\end{enumerate}

\section*{Long Form Response (20 marks)}
\textit{Answer each question with a clear and complete explanation.}
\begin{enumerate}[label=\arabic*.]
\setcounter{enumi}{10}
\item Discuss the factors that contribute to the differing relaxation times, T$_{1}$ and T$_{2}$, of a nucleus, with an emphasis on the influence of molecular motions at the Larmor frequency.
\item Discuss the concept of ionic strength and analyze why a 0.050 M solution of AlCl 3 exhibits higher ionic strength compared to other ionic solutions.
\end{enumerate}

\vfill
\noindent\textit{End of Paper}
\end{document}